\chapter{The bf16 Arm: Three Qwen Sizes and Llama~3.2}
\label{chap:s2-10}

\section{Goal}

The precision order set on 2026-11-14 is: every roster model in bf16 first, then every model in 8-bit,
4-bit last. This chapter reports the first models through the shared backend of
Chapter~\ref{chap:s2-9}: Qwen2.5-VL at three sizes (3B, 7B, 32B) with the full battery, and
Llama~3.2-11B-Vision verdict-only. The question for the Qwen arm is what model size changes when the
family, the prompt, the precision and the harness are all held fixed.

\section{What Was Done}

\textbf{Runs.} All four tiers ran on Scholar A40s into one results tree,
\texttt{results/scholar-a40-rf/}, with \texttt{XAI\_REGION\_FRAME=rescaled}, on the same 3,882
(image, rule) pairs as every earlier tier.
\begin{itemize}
\item \texttt{qwen3b-bf16}: full battery, one A40, 2.25~s per call.
\item \texttt{qwen7b-bf16}: the tier of Chapter~\ref{chap:s2-8}, unchanged; listed here for comparison.
Its Stage~2 also ran on the 7,457 compliant images outside the sample, so both 7B tiers now have a
verdict for every one of the 10,013 images.
\item \texttt{qwen32b-bf16}: full battery split over two A40s (33.2~GB on the first), 7.08~s per call.
Stages~4, 5, 6 and 8 took nearly three days of wall time, including a 19-hour stall (below), because
each job holds both GPUs the account may use and so they ran one at a time.
\item \texttt{llama32v-11b}: verdict-only chain (Stages 0--3, 6, 7), one A40, 3.65~s per call.
Its image tokens enter through cross-attention and form no grid over the image, so Stages~4, 5, 8, 9
and 10 do not apply (roster D4). It is asked for no region.
\end{itemize}

\textbf{Copy-back.} The tree was copied back as a tar, not through git, and all 191 data files were
checked by sha256 against Scholar. The Stage~5 attention grids (3,882 files per tier for 3B and 32B) are
committed with it. The 32B report regenerates locally byte for byte from the copied files. The verdict-only
chain does not run Stage~11, so the Llama report was produced locally by the same script.

\textbf{Environment B.} The parity run of \texttt{qwen7b-bf16} in environment B (\texttt{transformers}
5.20.2) on 96 rows agreed on the verdict in 94 (97.9\%), gave byte-identical replies in 48 (50.0\%), and
gave regions with median IoU 0.996 where both replies had one. This is the same order of disagreement as
two machines in environment A (96.9\% verdicts, 42.7\% replies; Chapter~\ref{chap:s2-1}). Models that
need environment B can be compared with the environment A tiers at the verdict level.

\section{Results}

\begin{table}[tbp]
\centering
\small
\begin{tabular}{lrrrr}
\toprule
\textbf{metric} & \textbf{3B} & \textbf{7B} & \textbf{32B} & \textbf{Llama 11B} \\
\midrule
\multicolumn{5}{l}{\emph{Stage 2 --- verdicts}} \\
sensitivity                         & 86.7 & 93.8 & 99.8 & 86.1 \\
specificity                         & 12.3 & 36.0 & 15.8 & 23.2 \\
YES share                           & 85.5 & 74.2 & 89.5 & 80.0 \\
\multicolumn{5}{l}{\emph{Stages 3, 9, 10 --- the stated explanation}} \\
no region stated                    & 30.3 & 20.5 & 8.8 & --- \\
explanation correct                 & 24.7 & 33.8 & 47.6 & --- \\
\quad IoU $\geq$ 0.30               & 48.5 & 55.8 & 65.2 & --- \\
\quad right answer, wrong reason    & 75.3 & 66.2 & 52.4 & --- \\
explanation supported (attention)   & 68.3 & 77.6 & 80.0 & --- \\
\multicolumn{5}{l}{\emph{Stages 4, 5 --- does the model use the region?}} \\
flip, region blacked out            & 14.4 & 27.9 & 15.7 & --- \\
\quad artefact floor (NO$\to$YES)   & 55.3 & 33.8 & 32.8$^{*}$ & --- \\
attention beats area null           & 97.9 & 97.6 & 87.8 & --- \\
\multicolumn{5}{l}{\emph{Stages 6, 8 --- stability and robustness}} \\
verdict unanimous over 5 reruns     & 10.4 & 47.0 & 72.9 & 40.6 \\
majority overturns greedy verdict   & 40.4 & 17.4 & 7.5 & 27.3 \\
verdict survives blur               & 62.0 & 68.8 & 73.5 & --- \\
verdict survives occlusion          & 69.1 & 79.4 & 83.2 & --- \\
\bottomrule
\end{tabular}
\caption{The bf16 arm, all rates in percent. Bootstrap 95\% intervals are in the per-tier reports
(\texttt{results/scholar-a40-rf/report\_<tier>.txt}); intervals on the Stage~2 and Stage~10 rows are at
most $\pm$3.5 points. $^{*}$ $n=64$: the 32B model answers NO to so few rows that its floor is loose
(95\% CI 21.9--45.3). Stage~8 samples per tier from the tier's own verdicts, so its $n$ differs (3B 1,288;
7B 1,599; 32B 1,206); every planned call ran.}
\label{tab:s210}
\end{table}

\textbf{Specificity does not improve with size.} Sensitivity rises with size
(86.7, 93.8, 99.8\%) but specificity peaks at 7B (12.3, 36.0, 15.8\%). The 32B model answers YES to 89.5\%
of all questions and misses 3 of 1,327 hazard pairs. A higher sensitivity bought by saying YES more often is
not better detection, so the false-alarm rate stays a headline number for every tier (arch doc §5.3).
Llama~3.2 is level with the 3B model on sensitivity (86.1\%) and between 3B and 7B on specificity
(23.2\%).

\textbf{The stated explanation improves steadily with size.} Explanation correct rises from 24.7 to 33.8
to 47.6\%, and every component moves the same way: the region is stated more often, overlaps the annotated
box more often, and is more often backed by attention (Stage~9). ``Right answer, wrong reason'' falls from
three in four to one in two. As in every earlier tier, explanations are more often correct on poor-info
images than on rich-info ones (32B: 53.6 vs.\ 43.6\%).

\textbf{Black masking says little at any size.} In no bf16 tier does the pooled flip rate clear its
artefact floor: blacking out a patch on an image the model called compliant turns the verdict to YES at
least as often as blacking out the stated region turns YES to NO. For 32B the flip rate is half the floor.
Rule~4 is the one rule whose flip rate cleared the floor at 7B (54.6\%); at 32B it is level with the floor
(32.0\%). Stage~4 cannot separate ``the model uses its region'' from ``the model ignores edits'' for
these tiers (F4).

\textbf{Attention lands in the stated region at every size, with one exception.} The 3B and 7B models beat
the area null on 97--98\% of rows. The 32B model beats it on 87.8\%, and the whole drop is rule~3 (59.6\%,
against 92.5\% at 7B). Rules~1, 2 and 4 stay at 99\%. This is the first rule-specific attention failure in
the study and is not yet explained; a look at the rule~3 grids is the next check.

\textbf{Larger models are far more stable.} At temperature 0.7, the five reruns all agree on 10.4\% of rows
for 3B, 47.0\% for 7B and 72.9\% for 32B, and the majority overturns the greedy verdict on 40.4, 17.4 and
7.5\%. The 3B model is unstable enough that its greedy verdicts should be read as one draw, not as the
model's answer. Robustness to blur and occlusion also improves with size, by about 10 points from 3B to
32B.

\section{Problems and Solutions}

\begin{itemize}
\item \textbf{The 32B Stage~6 latency guard tripped on a healthy run.} The guard compared a 25~s ceiling
with the time per row, and a Stage~6 row is five calls (36.5~s at 7.3~s per call). The chains stopped for
about 19 hours. The guard now reads each call's own time (\texttt{reply\_latency\_s}, commit
\texttt{4ba671f}); Stage~8 got the same fix. Stage~4 still sums its two calls, which is safe at 14~s for
32B but may trip on slower 8-bit tiers.
\item \textbf{The 3B Stage~8 parse guard tripped on 1.4\% of replies.} The guard was per row, so one
unreadable reply among six failed the row. It is now per reply, as Stage~6 has been since
Chapter~\ref{chap:s2-5} (commit \texttt{d40875b}).
\item \textbf{Llama sampled with its own top-$p$.} The checkpoint's generation config sets
\texttt{top\_p=0.9}, so Llama's Stage~6 was not sampling the same way as Qwen's. Stage~6 now pins
\texttt{top\_p=1.0} and \texttt{top\_k=50} for every tier (Qwen unchanged), and Llama's Stage~6 was rerun.
The first file is kept in \texttt{superseded/}.
\item \textbf{A two-GPU job never starts behind one-GPU jobs.} The account may hold two GPUs. While
one-GPU jobs kept queueing, the 32B sizing job waited indefinitely; pending one-GPU jobs were held behind
it with a dependency.
\end{itemize}

\section{Standard vs.\ Non-Standard Choices}

\textbf{Standard.} Comparing sizes within one model family at one precision is how scaling is usually
measured. Multi-GPU sharding of the 32B model (\texttt{device\_map=auto}) is the stock \texttt{accelerate}
path.

\textbf{Non-standard, and stated in the paper.}
\begin{itemize}
\item Explanation metrics are computed only on the rows where a model gives the explanation (Stage~10 on
hazard pairs answered YES; Stages~4, 5, 9 on rows with a region). A model with a stronger YES bias is
scored on more rows. The $n$ is reported beside every rate for this reason.
\item Qwen keeps its checkpoint's \texttt{repetition\_penalty} of 1.05 under greedy decoding; Llama's is
1.0. This is each model's shipped default and is not pinned; it is documented, not equalised.
\end{itemize}

\section{Implications}

\begin{itemize}
\item Within one family, size buys explanation quality and stability, not specificity. A paper claim of
the form ``larger models explain better'' holds here; ``larger models detect better'' does not.
\item The 32B tier gives the clearest case of a model that is right for the stated reason most of the time
(47.6\% correct, 80\% attention-supported) while flagging 84\% of compliant images. Explanation quality and
verdict accuracy have to be reported side by side.
\item Next: the rule~3 attention grids at 32B, then the remaining roster models in bf16 (adapters and
Stage~0 sizing per model), then the 8-bit arm.
\end{itemize}
